The 8th IOAA organizers plan to send to the \textbf{antipode} the official
flag using a ballistic projectile. The projectile will be launched from
Romania, and the rotation of the Earth will be neglected.

\begin{description}
\item[(a)] Calculate the coordinates of the target-point if the launch-point
  coordinates are: $\varphi_{\mathrm{Romania}} = 44^\circ$ North;
  $\lambda_{\mathrm{Romania}} = 30^\circ$ East.
\item[(b)] Determine the elements of the launching-speed vector, regardless to
  the center of the Earth, in order that the projectile should hit the target.
\item[(c)] Calculate the velocity of the projectile when it hits the target.
\item[(d)] Calculate the minimum velocity of the projectile.
\item[(e)] Calculate the flying-time of the projectile, from the launch-moment
  to the impact one. You will use the value of the gravitational acceleration
  at Earth surface $g_0 = 9.81$\,$\mathrm{m\,s^{-2}}$; the Earth radius
  $R = 6370$\,km.
\item[(f)] Evaluate the possibility that the projectile to be seen with the
  naked eye in the moment that it passes at the maximum distance from the
  Earth. You will use the following values: the Moon albedo
  $\alpha_M = 0.12$; the Moon radius $R_M = 1738$\,km; the Earth--Moon
  distance $r_{EL} = 384\,400$\,km; the apparent magnitude of the full moon
  $m_M = -12.7^{\mathrm{m}}$. You assume that the projectile is perfectly
  metallic sphere with radius
  $r_{\mathrm{projectile}} = 400 \times 10^{-3}$\,m and with perfectly
  reflective surface.
\end{description}
